\documentclass[11pt]{scrartcl}

\usepackage[T1]{fontenc}
\usepackage[utf8]{inputenc}
\usepackage{lmodern}
\usepackage{amsmath,amssymb,amsthm,amsfonts,mathtools}
\usepackage{geometry}
\geometry{margin=1.1in}
\usepackage{microtype}
\usepackage{enumitem}
\usepackage[numbers]{natbib}
\usepackage[hidelinks]{hyperref}

\numberwithin{equation}{section}

\newtheorem{theorem}{Theorem}[section]
\newtheorem{proposition}[theorem]{Proposition}
\newtheorem{lemma}[theorem]{Lemma}
\newtheorem{corollary}[theorem]{Corollary}
\theoremstyle{definition}
\newtheorem{remark}[theorem]{Remark}

\newcommand{\N}{\mathbb N}
\newcommand{\Z}{\mathbb Z}
\newcommand{\R}{\mathbb R}
\newcommand{\C}{\mathbb C}
\newcommand{\T}{\mathbb T}
\newcommand{\1}{\mathbf 1}
\newcommand{\abs}[1]{\left|#1\right|}
\newcommand{\norm}[1]{\left\|#1\right\|}
\newcommand{\ee}[1]{e\!\left(#1\right)}

\title{Fixed-shift Liouville pair correlations: exact reductions, projected transport, and local quadratic obstructions}
\author{Míngshū Wáng\thanks{Corresponding author. Email: adriel.datstat@gmail.com}}
\date{April 2026}

\begin{document}
\maketitle

\begin{abstract}
For a fixed nonzero shift $h$, we study the Liouville pair-correlation sums
\[
A_h(N):=\sum_{n\le N}\lambda(n)\lambda(n+h).
\]
We develop a critical-scale reduction theory for the asymptotic formula $A_h(N)=o(N)$, the fixed-shift two-point case of Chowla's conjecture, at the scale $H=X^{1/2}$. Passing to the derived sequence $a_h(n)=\lambda(n)\lambda(n+h)$, we separate the direct absolute first-moment implication from the signed Fej\'er-kernel and dyadic-shell identities, and we record exact quartic Ces\`aro, Fourier--Dirichlet, dyadic-shell, and first-moment reformulations of the critical short-block criterion. We then prove the correct fixed-shift local $U^2$ reduction: the normalized local quantity is a fourth-power quantity, so block sums are controlled through a fourth-power estimate and a Young/H\"older averaging step, not by a false linear pointwise bound. Finally we introduce a projected packet formalism for bad blocks. On the projected side we establish an exact affine-shear identity, a rigidity theorem, and an exact transport-defect identity, and we show that a large projected Fej\'er/shear packet forces a local quadratic-phase obstruction for $\lambda$. The output is a structured conditional reduction of the remaining obstruction to projected bias, transport rigidity, and projected long-range packet formation.
\end{abstract}

\section{Introduction and statement of results}

Let $\lambda(n)=(-1)^{\Omega(n)}$ be the Liouville function. For a fixed nonzero shift $h\in\Z$, we consider the pair-correlation sum
\[
A_h(N):=\sum_{n\le N}\lambda(n)\lambda(n+h).
\]
The asymptotic formula
\[
A_h(N)=o(N)
\qquad (N\to\infty)
\]
is the fixed-shift two-point case of Chowla's conjecture \cite{chowla1965riemann}. Major progress on correlations of multiplicative functions is now available in logarithmically averaged settings, at almost all scales, and through short-interval, Fourier-uniformity, and higher-uniformity theorems on average \cite{tao2016logarithmically,taoteravainen2019almostall,matomaki2016multiplicative,matomaki2020fourier,matomaki2023higher,greentao2012mobiusnil,matomakiShao2021poly}. The fixed-shift pointwise problem nevertheless remains especially rigid. The purpose of this paper is to isolate that rigidity at the critical scale
\[
H=X^{1/2},
\]
and to show that the remaining obstruction can be expressed in a sharply structured local form.

Our analysis is carried out entirely at one fixed shift. Starting from the critical short-block mean-square criterion of Section~\ref{sec:short-block}, we pass to the derived sequence
\[
a_h(n):=\lambda(n)\lambda(n+h),
\]
and develop a chain of exact reformulations and reductions for its short sums and associated energy profile. In particular, we obtain quartic Ces\`aro, Fourier--Dirichlet, dyadic-shell, and short-interval first-moment criteria, together with a positive multiscale absorption theorem. The Fourier-kernel side of the argument is classical in spirit \cite{zygmund1959trigonometric,grafakos2008classical}, while the multiscale organization is deterministic and chaining-style rather than generic chaining in Talagrand's stochastic-process sense \cite{talagrand2005generic}. We then record a fixed-shift local $U^2$ reduction in the fourth-power normalization used in the formal development. This point is important: the averaged local $U^2$ input is consumed through a fourth-power estimate and a scalar Young inequality, not through a linear pointwise bound in the normalized fourth-power mass. Finally we introduce a projected packet formalism for bad blocks. On the projected side we establish an exact affine-shear identity, a rigidity theorem, and an exact transport-defect identity, and we show that a large projected Fej\'er/shear packet forces a local quadratic-phase obstruction for $\lambda$, in the broader context of recent polynomial-phase and higher-uniformity results for multiplicative functions in short intervals \cite{matomaki2023higher,greentao2012mobiusnil,matomakiShao2021poly}. In this way the fixed-shift problem is reduced to a precise final dichotomy between projected bias and projected long-range packet formation.

A basic global consequence of the short-block framework is the following conditional implication. It is stated in absolute first-moment form; this is stronger information than signed Fej\'er cancellation, because it controls the total mass of bad short blocks rather than cancellations between shells.

\begin{theorem}[Conditional fixed-shift theorem]\label{thm:main}
Fix $h\neq 0$. Assume that, with $H=\lfloor X^{1/2}\rfloor$, one has
\begin{equation}\label{eq:main-first-moment-hyp}
\sup_{1\le M\le H}
\frac{1}{XM}
\sum_{X\le x\le 2X-M}
\abs{\sum_{x<n\le x+M}\lambda(n)\lambda(n+h)}
\to 0
\qquad (X\to\infty).
\end{equation}
Then
\[
A_h(N)=o(N)
\qquad (N\to\infty).
\]
\end{theorem}

Section~\ref{sec:first-moment} contains the direct first-moment implication and its dyadic-shell consequence. Section~\ref{sec:local-u2} records the corrected fourth-power local $U^2$ reduction. Section~\ref{sec:projected-packets} develops the projected transport and packet framework that captures the final transition in the reduction.

\section{The critical short-block criterion}
\label{sec:short-block}

Fix $h\neq 0$ and define
\[
a_h(n):=\lambda(n)\lambda(n+h),
\qquad
A_h(N):=\sum_{n\le N}a_h(n),
\]
and for integers $x$ and $H\ge 1$ define
\[
B_h(x;H):=\sum_{x<n\le x+H}a_h(n)=A_h(x+H)-A_h(x).
\]

\begin{theorem}[Critical short-block criterion]
\label{thm:short-block-criterion}
Fix $h\neq 0$. Assume that for $H=\lfloor X^{1/2}\rfloor$ one has
\begin{equation}
\label{eq:short-block-criterion-hyp}
\sum_{X\le x\le 2X-H}\abs{B_h(x;H)}^2=o(XH^2)
\qquad (X\to\infty).
\end{equation}
Then
\begin{equation}
\label{eq:short-block-criterion-conc}
A_h(N)=o(N)
\qquad (N\to\infty).
\end{equation}
\end{theorem}

\begin{proof}
Write
\[
\eta(X):=\frac{1}{XH^2}\sum_{X\le x\le 2X-H}\abs{B_h(x;H)}^2,
\]
so that $\eta(X)\to 0$ by \eqref{eq:short-block-criterion-hyp}. Fix a parameter $0<\tau\le 1$, to be chosen later, and define the bad set
\[
\mathcal E_\tau:=\{x\in [X,2X-H]\cap\Z:\ \abs{B_h(x;H)}>\tau H\}.
\]
By Markov's inequality,
\[
\#\mathcal E_\tau\cdot \tau^2 H^2
\le
\sum_{X\le x\le 2X-H}\abs{B_h(x;H)}^2
=\eta(X)XH^2,
\]
so
\begin{equation}
\label{eq:bad-set-bound}
\#\mathcal E_\tau\le \eta(X)\tau^{-2}X.
\end{equation}

Now split starting points by residue class modulo $H$. For $r\in\{0,1,\dots,H-1\}$ set
\[
x_{r,j}:=X+r+jH.
\]
Averaging over $r$ shows that there exists a residue class $r=r_X$ such that
\[
\#\{j:\ x_{r,j}\in \mathcal E_\tau\}
\le
\frac{\#\mathcal E_\tau}{H}
\le
\eta(X)\tau^{-2}\frac{X}{H}.
\]
Fix such an $r$ and write $x_j:=x_{r,j}$.

Let $N\in[X,2X]$, and choose $m$ so that
\[
x_m\le N < x_{m+1},
\]
with the convention $m=-1$ if $N<x_0$. Then
\[
A_h(N)-A_h(x_0)
=
\sum_{j=0}^{m-1}\bigl(A_h(x_{j+1})-A_h(x_j)\bigr)+O(H)
=
\sum_{j=0}^{m-1} B_h(x_j;H)+O(H),
\]
since the tail interval between $x_m$ and $N$ has length $<H$ and $\abs{a_h(n)}\le 1$.

Among the blocks $x_j$, every good block contributes at most $\tau H$ in absolute value, while the total number of bad blocks is at most $\eta(X)\tau^{-2}X/H$. Since $mH\ll X$, we obtain
\[
\abs{A_h(N)-A_h(x_0)}
\ll
\tau X + \eta(X)\tau^{-2}X + H.
\]
Also $x_0=X+r$ with $0\le r<H$, so
\[
\abs{A_h(x_0)-A_h(X)}\le H.
\]
Hence
\begin{equation}
\label{eq:blockwise-prefix-bound}
\sup_{X\le N\le 2X}\abs{A_h(N)-A_h(X)}
\ll
\tau X + \eta(X)\tau^{-2}X + H.
\end{equation}
Choose $\tau:=\eta(X)^{1/3}$. Then
\[
\tau X + \eta(X)\tau^{-2}X \ll \eta(X)^{1/3}X,
\]
and because $H=X^{1/2}+O(1)=o(X)$, \eqref{eq:blockwise-prefix-bound} becomes
\begin{equation}
\label{eq:blockwise-prefix-oX}
\sup_{X\le N\le 2X}\abs{A_h(N)-A_h(X)}=o(X).
\end{equation}

Now let $X_j:=2^j$. For any $N$ with $2^K\le N<2^{K+1}$,
\[
A_h(N)
=
A_h(1)+\sum_{j=0}^{K-1}\bigl(A_h(2^{j+1})-A_h(2^j)\bigr)+\bigl(A_h(N)-A_h(2^K)\bigr).
\]
Applying \eqref{eq:blockwise-prefix-oX} on each dyadic interval $[2^j,2^{j+1}]$ shows that
\[
A_h(2^{j+1})-A_h(2^j)=o(2^j),
\qquad
A_h(N)-A_h(2^K)=o(2^K).
\]
Therefore, for every $\varepsilon>0$, all sufficiently large $j$ satisfy
\[
\abs{A_h(2^{j+1})-A_h(2^j)}\le \varepsilon 2^j,
\qquad
\abs{A_h(N)-A_h(2^K)}\le \varepsilon 2^K,
\]
so
\[
\abs{A_h(N)}\le O_\varepsilon(1)+\varepsilon\sum_{j=0}^{K}2^j\ll O_\varepsilon(1)+\varepsilon N.
\]
Dividing by $N$ and letting $N\to\infty$, then $\varepsilon\to 0$, gives \eqref{eq:short-block-criterion-conc}.
\end{proof}

\section{Energy profile and quartic correlations}
\label{sec:energy-profile}

We now pass from the short-block criterion to the exact quadratic and quartic identities attached to the derived sequence $a_h$. Throughout this section, let $X\ge 1$ be large and set
\[
H:=\lfloor X^{1/2}\rfloor.
\]
For integers $M\ge 1$, define the short-block energy
\[
E_h(M;X):=\sum_{X\le x\le 2X-M}\abs{B_h(x;M)}^2.
\]
Also define the quartic correlation sums
\[
P_h(t;X):=\sum_{X<n\le 2X-t}a_h(n)a_h(n+t)
\qquad (t\ge 1),
\]
and the Ces\`aro partial sums
\[
C_h(M;X):=\sum_{1\le t\le M}P_h(t;X).
\]

\begin{proposition}[Exact energy profile]
\label{prop:energy-profile}
For every integer $M\ge 1$,
\begin{equation}
\label{eq:energy-profile}
E_h(M;X)=MX+2\sum_{1\le t<M}(M-t)P_h(t;X)+O_h(M^3).
\end{equation}
\end{proposition}

\begin{proof}
Let
\[
a_{h,X}(n):=a_h(n)\,\mathbf 1_{X<n\le 2X}.
\]
Set
\[
\widetilde B_h(x;M):=\sum_{x<n\le x+M}a_{h,X}(n).
\]
For $X\le x\le 2X-M$ one has $\widetilde B_h(x;M)=B_h(x;M)$. Outside this range there are only $O(M)$ values of $x$ for which $\widetilde B_h(x;M)\neq 0$, and trivially $\abs{\widetilde B_h(x;M)}\le M$. Hence
\begin{equation}
\label{eq:full-vs-truncated}
\sum_{x\in\Z}\abs{\widetilde B_h(x;M)}^2=E_h(M;X)+O(M^3).
\end{equation}

Expanding the square gives the exact block-energy identity
\[
\sum_{x\in\Z}\abs{\widetilde B_h(x;M)}^2
=
\sum_{|t|<M}(M-|t|)\sum_{n\in\Z}a_{h,X}(n)a_{h,X}(n+t).
\]
Since $X>|h|$ for $X$ large, we have $a_h(n)^2=1$ on $(X,2X]$, and so
\[
\sum_{n\in\Z}a_{h,X}(n)^2=X+O_h(1).
\]
For $t\ge 1$,
\[
\sum_{n\in\Z}a_{h,X}(n)a_{h,X}(n+t)=P_h(t;X),
\]
and the contribution of negative $t$ is the same by symmetry. Therefore
\[
\sum_{x\in\Z}\abs{\widetilde B_h(x;M)}^2
=MX+2\sum_{1\le t<M}(M-t)P_h(t;X)+O_h(M).
\]
Combining this with \eqref{eq:full-vs-truncated} proves \eqref{eq:energy-profile}.
\end{proof}

\begin{corollary}[Discrete derivative of the energy profile]
\label{cor:discrete-derivative}
For every integer $1\le M\le H$,
\begin{equation}
\label{eq:discrete-derivative}
E_h(M+1;X)-E_h(M;X)=X+2C_h(M;X)+O_h(M^2).
\end{equation}
\end{corollary}

\begin{proof}
Subtract \eqref{eq:energy-profile} at $M+1$ and $M$. The diagonal term contributes $X$, while the quartic terms telescope to $2\sum_{1\le t\le M}P_h(t;X)=2C_h(M;X)$. The error term changes by $O_h(M^2)$.
\end{proof}

\begin{corollary}[A uniform slope criterion]
\label{cor:uniform-slope-criterion}
If
\[
\sup_{1\le M\le H}\abs{E_h(M+1;X)-E_h(M;X)-X}=o(XH),
\]
then \eqref{eq:short-block-criterion-conc} holds.
\end{corollary}

\begin{proof}
By \eqref{eq:discrete-derivative},
\[
\sup_{1\le M\le H}\abs{C_h(M;X)}=o(XH)
\]
(because $M^2\ll H^2\asymp X=o(XH)$). Then
\[
\sum_{1\le t<H}(H-t)P_h(t;X)=\sum_{M=1}^{H-1}C_h(M;X)=o(XH^2),
\]
and now \eqref{eq:energy-profile} at $M=H$ yields
\[
E_h(H;X)=HX+o(XH^2)+O(H^3)=o(XH^2),
\]
since $HX\asymp X^{3/2}=o(X^2)=o(XH^2)$ and $H^3\asymp X^{3/2}=o(XH^2)$. The conclusion follows from Theorem~\ref{thm:short-block-criterion}.
\end{proof}


\begin{remark}[Signed Fej\'er cancellation and Abel summation]
\label{rem:signed-fejer-abel}
The preceding identities also explain a useful boundary of the method. If
\[
S_h(M;X):=\sum_{1\le t\le M}P_h(t;X),
\]
then the Fej\'er-weighted off-diagonal term satisfies the exact Abel identity
\[
\sum_{1\le t<H}(H-t)P_h(t;X)
=\sum_{M=1}^{H-1}S_h(M;X).
\]
Consequently a signed hypothesis such as
\[
\sup_{1\le M\le H}\abs{S_h(M;X)}=o(XH)
\]
would imply the critical energy estimate by a direct summation-by-parts argument, together with Proposition~\ref{prop:energy-profile}. This signed route is deliberately not the same as the absolute first-moment hypothesis used in Theorem~\ref{thm:main}: signed Fej\'er cancellation can hide cancellation between positive and negative shell contributions, whereas the absolute first moment controls the total mass of large short blocks and feeds into the bad-block and packet reductions below.
\end{remark}

\section{Fourier reformulation and positive multiscale absorption}
\label{sec:fourier}

The discrete derivative identity naturally leads to a Fourier reformulation. The short-block energies are represented by Fej\'er kernels, while their discrete derivatives are represented by Dirichlet kernels; this is standard Fourier analysis on the torus \cite{zygmund1959trigonometric,grafakos2008classical}. This is the point at which the oscillatory obstruction is converted into a positive multiscale quantity.

Define
\[
A_{h,X}(\alpha):=\sum_{X<n\le 2X}a_h(n)\ee{\alpha n},
\qquad \alpha\in \R/\Z.
\]
For $M\ge 1$, let
\[
F_M(\alpha):=\left|\sum_{m=1}^M \ee{\alpha m}\right|^2,
\qquad
D_M(\alpha):=\sum_{|t|\le M}\ee{\alpha t}.
\]

\begin{proposition}[Fej\'er representation]
\label{prop:fejer-representation}
For every integer $M\ge 1$,
\begin{equation}
\label{eq:fejer-representation}
E_h(M;X)=\int_0^1 \abs{A_{h,X}(\alpha)}^2F_M(\alpha)\,d\alpha+O_h(M^3).
\end{equation}
Moreover, for $1\le M\le H$,
\begin{equation}
\label{eq:dirichlet-increment}
E_h(M+1;X)-E_h(M;X)=\int_0^1 \abs{A_{h,X}(\alpha)}^2D_M(\alpha)\,d\alpha+O_h(M^2).
\end{equation}
\end{proposition}

\begin{proof}
Let $K_M(n):=\mathbf 1_{1\le n\le M}$. Then
\[
\widetilde B_h(x;M)=(a_{h,X}*K_M)(x),
\]
where $a_{h,X}$ is as in the proof of Proposition~\ref{prop:energy-profile}. By Parseval on $\Z$,
\[
\sum_{x\in\Z}\abs{\widetilde B_h(x;M)}^2
=\int_0^1\abs{A_{h,X}(\alpha)}^2\abs{\widehat K_M(\alpha)}^2\,d\alpha
=\int_0^1\abs{A_{h,X}(\alpha)}^2F_M(\alpha)\,d\alpha.
\]
Using \eqref{eq:full-vs-truncated} proves \eqref{eq:fejer-representation}. The identity
\[
F_{M+1}(\alpha)-F_M(\alpha)=D_M(\alpha)
\]
follows by comparing Fourier coefficients, and then \eqref{eq:dirichlet-increment} follows by subtraction.
\end{proof}

\begin{lemma}[A positive multiscale domination]
\label{lem:positive-domination}
For every integer $H\ge 1$ and every $\alpha\in\R/\Z$,
\begin{equation}
\label{eq:positive-domination}
\sup_{1\le L\le H}\abs{D_L(\alpha)-1}
\ll
\sum_{M=1}^H \frac{F_M(\alpha)}{M^2}.
\end{equation}
\end{lemma}

\begin{proof}
Let $\delta:=\norm{\alpha}$, where $\norm{\alpha}$ denotes the distance from $\alpha$ to the nearest integer. The standard Dirichlet-kernel bound gives
\[
\sup_{1\le L\le H}\abs{D_L(\alpha)-1}\ll \min(H,\delta^{-1}).
\]
It therefore suffices to show
\[
\sum_{M=1}^H \frac{F_M(\alpha)}{M^2}\gg \min(H,\delta^{-1}).
\]
If $\delta\le 1/(2H)$, then $\delta\le 1/(2M)$ for every $1\le M\le H$, and the usual estimate
\[
F_M(\alpha)=\left|\frac{\sin(\pi M\alpha)}{\sin(\pi\alpha)}\right|^2\gg M^2
\]
shows that the sum is $\gg H$.

If $\delta>1/(2H)$, let $M_0:=\lfloor (2\delta)^{-1}\rfloor$. Then $1\le M_0\le H$ and still $\delta\le 1/(2M)$ for every $1\le M\le M_0$, so $F_M(\alpha)\gg M^2$ throughout that range. Hence
\[
\sum_{M=1}^H \frac{F_M(\alpha)}{M^2}\gg M_0\asymp \delta^{-1}.
\]
This proves \eqref{eq:positive-domination}.
\end{proof}

\begin{theorem}[Positive multiscale criterion]
\label{thm:positive-multiscale-criterion}
Assume that
\begin{equation}
\label{eq:positive-multiscale-hyp}
\sum_{M=1}^{H}\frac{E_h(M;X)}{M^2}=o(XH)
\qquad (X\to\infty).
\end{equation}
Then \eqref{eq:short-block-criterion-conc} holds.
\end{theorem}

\begin{proof}
By Proposition~\ref{prop:fejer-representation},
\[
\int_0^1 \abs{A_{h,X}(\alpha)}^2\,d\alpha
=\sum_{X<n\le 2X}a_h(n)^2
=X+O_h(1).
\]
Hence, by Corollary~\ref{cor:uniform-slope-criterion}, it is enough to show that
\[
\sup_{1\le L\le H}\left|\int_0^1 \abs{A_{h,X}(\alpha)}^2(D_L(\alpha)-1)\,d\alpha\right|=o(XH).
\]
Applying Lemma~\ref{lem:positive-domination}, then interchanging sum and integral, and finally using Proposition~\ref{prop:fejer-representation}, we obtain
\[
\sup_{1\le L\le H}\left|\int_0^1 \abs{A_{h,X}(\alpha)}^2(D_L(\alpha)-1)\,d\alpha\right|
\ll
\sum_{M=1}^{H}\frac{1}{M^2}\int_0^1 \abs{A_{h,X}(\alpha)}^2F_M(\alpha)\,d\alpha
\]
\[
=
\sum_{M=1}^{H}\frac{E_h(M;X)}{M^2}+O\!\left(\sum_{M=1}^{H}M\right)
=
\sum_{M=1}^{H}\frac{E_h(M;X)}{M^2}+O(H^2).
\]
Since $H^2\asymp X=o(XH)$, hypothesis \eqref{eq:positive-multiscale-hyp} yields the desired bound.
\end{proof}

\section{Dyadic shells, tree weights, and harmonic weights}
\label{sec:dyadic-shells}

We next reorganize the multiscale sum shell-by-shell in the shift variable $t$. This isolates the precise dyadic local cancellation required of the quartic family.

Let $J\ge 0$ be defined by
\[
2^{J-1}<H\le 2^J.
\]
For $0\le k\le J-1$, write
\[
I_k:=\{t\in\Z: 2^k\le t<2^{k+1}\}.
\]
For integers $u\in I_k$, define the shell partial sums
\[
Q_{h,k}(u;X):=\sum_{2^k\le t\le u} P_h(t;X).
\]
Also define the tree-stopped multiscale energy
\[
\mathcal T_h(X):=\sum_{j=0}^{J}\frac{E_h(2^j;X)}{2^j}.
\]

\begin{proposition}[Tree weight identity]
\label{prop:tree-weight-identity}
One has
\begin{equation}
\label{eq:tree-weight-identity}
\mathcal T_h(X)=X(J+1)+2\sum_{1\le t<2^J}\omega_J(t)P_h(t;X)+O_h(H^2),
\end{equation}
where
\[
\omega_J(t):=\sum_{\substack{0\le j\le J \\ 2^j>t}}\left(1-\frac{t}{2^j}\right).
\]
If $t\in I_k$, then
\begin{equation}
\label{eq:omega-affine}
\omega_J(t)=A_k-B_kt,
\qquad
A_k:=J-k,
\qquad
B_k:=\sum_{j=k+1}^{J}2^{-j}.
\end{equation}
In particular,
\begin{equation}
\label{eq:omega-size}
A_k\ll J-k+1,
\qquad
B_k\asymp 2^{-k}.
\end{equation}
\end{proposition}

\begin{proof}
Insert \eqref{eq:energy-profile} with $M=2^j$ into the definition of $\mathcal T_h(X)$ and interchange the order of summation. For $t\in I_k$, the condition $2^j>t$ is equivalent to $j\ge k+1$, which gives the affine formula \eqref{eq:omega-affine}. The size bounds are immediate.
\end{proof}

\begin{proposition}[Harmonic weight identity]
\label{prop:harmonic-weight-identity}
Define
\[
\Sigma_h(X):=\sum_{M=1}^{H}\frac{E_h(M;X)}{M^2}.
\]
Then
\begin{equation}
\label{eq:harmonic-weight-identity}
\Sigma_h(X)=X\sum_{M=1}^{H}\frac1M+2\sum_{1\le t<H}\nu_H(t)P_h(t;X)+O_h(H^2),
\end{equation}
where
\[
\nu_H(t):=\sum_{M=t+1}^{H}\frac{M-t}{M^2}
=\sum_{M=t+1}^{H}\frac1M-t\sum_{M=t+1}^{H}\frac1{M^2}.
\]
Moreover, if $t\in I_k$, then
\begin{equation}
\label{eq:nu-size}
\nu_H(t)\ll J-k+1,
\end{equation}
and for $2^k\le u<2^{k+1}-1$,
\begin{equation}
\label{eq:nu-difference}
\nu_H(u)-\nu_H(u+1)=\sum_{M=u+1}^{H}\frac1{M^2}\asymp 2^{-k}.
\end{equation}
\end{proposition}

\begin{proof}
Insert \eqref{eq:energy-profile} into the definition of $\Sigma_h(X)$ and interchange the order of summation:
\[
\Sigma_h(X)
=
X\sum_{M=1}^{H}\frac1M+2\sum_{1\le t<H}\left(\sum_{M=t+1}^{H}\frac{M-t}{M^2}\right)P_h(t;X)+O_h\!\left(\sum_{M=1}^{H}M\right).
\]
Since $\sum_{M=1}^{H}M\ll H^2$, this is \eqref{eq:harmonic-weight-identity}.

If $t\in I_k$, then
\[
\nu_H(t)\le \sum_{M=t+1}^{H}\frac1M \ll \log\!\frac{H}{t+1}+1\ll J-k+1,
\]
which gives \eqref{eq:nu-size}. Finally,
\[
\nu_H(u)-\nu_H(u+1)
=
\frac1{(u+1)^2}+\sum_{M=u+2}^{H}\frac1{M^2}
=
\sum_{M=u+1}^{H}\frac1{M^2},
\]
and since $u\asymp 2^k$ on $I_k$, this is $\asymp 2^{-k}$. This proves \eqref{eq:nu-difference}.
\end{proof}

\begin{theorem}[Dyadic local quartic criterion]
\label{thm:dyadic-local-quartic-criterion}
Assume that
\begin{equation}
\label{eq:dyadic-local-quartic-hyp}
\sup_{\substack{0\le k\le J-1 \\ u\in I_k}}\frac{\abs{Q_{h,k}(u;X)}}{X2^k}\to 0
\qquad (X\to\infty).
\end{equation}
Then \eqref{eq:positive-multiscale-hyp} holds, hence \eqref{eq:short-block-criterion-conc} holds.
\end{theorem}

\begin{proof}
We apply shellwise summation by parts to the weighted quartic sum in \eqref{eq:harmonic-weight-identity}. For a fixed shell $I_k=[2^k,2^{k+1}-1]$, write
\[
\Sigma_k:=\sum_{t\in I_k}\nu_H(t)P_h(t;X).
\]
Then
\[
\Sigma_k
=
\nu_H(2^{k+1}-1)Q_{h,k}(2^{k+1}-1;X)
+
\sum_{u=2^k}^{2^{k+1}-2}\bigl(\nu_H(u)-\nu_H(u+1)\bigr)Q_{h,k}(u;X).
\]
Using \eqref{eq:nu-size}, \eqref{eq:nu-difference}, and hypothesis \eqref{eq:dyadic-local-quartic-hyp}, we obtain
\[
\abs{\Sigma_k}
\ll
(J-k+1)\sup_{u\in I_k}\abs{Q_{h,k}(u;X)}
+
2^{-k}\sum_{u=2^k}^{2^{k+1}-2}\sup_{v\in I_k}\abs{Q_{h,k}(v;X)}
\]
\[
\ll
(J-k+1)\sup_{u\in I_k}\abs{Q_{h,k}(u;X)}
=
 o\!\bigl(X(J-k+1)2^k\bigr).
\]
Summing over $k$ and using
\[
\sum_{k=0}^{J-1}(J-k+1)2^k\ll 2^J\asymp H,
\]
we deduce
\[
\sum_{1\le t<H}\nu_H(t)P_h(t;X)=o(XH).
\]
Also
\[
X\sum_{M=1}^{H}\frac1M\ll X\log H=o(XH),
\qquad
H^2=o(XH).
\]
Hence \eqref{eq:harmonic-weight-identity} gives
\[
\Sigma_h(X)=o(XH),
\]
which is exactly \eqref{eq:positive-multiscale-hyp}. The conclusion now follows from Theorem~\ref{thm:positive-multiscale-criterion}.
\end{proof}

\begin{remark}
The dyadic local criterion isolates a scale-by-scale version of the quartic cancellation problem. The input is local in the $t$-variable and is therefore naturally compatible with shellwise stopping and summation-by-parts arguments.
\end{remark}

\section{A short-interval first-moment theorem}
\label{sec:first-moment}

We now isolate the role of the absolute first moment. There is a direct implication from
\[
\sup_{1\le M\le H}
\frac{1}{XM}\sum_{X\le x\le 2X-M}\abs{B_h(x;M)}=o(1)
\]
to the critical short-block criterion: at the single scale $M=H$, the pointwise bound $\abs{B_h(x;H)}\le H$ gives
\[
\sum_{X\le x\le 2X-H}\abs{B_h(x;H)}^2
\le
H\sum_{X\le x\le 2X-H}\abs{B_h(x;H)}=o(XH^2).
\]
Thus the first-moment implication itself is elementary once absolute values are available. We nevertheless keep the dyadic-shell bridge because it identifies exactly which local quartic shell quantities are controlled by the same absolute first moments, and these shell quantities feed the packet analysis below.

\begin{proposition}[Shells are dominated by first moments]
\label{prop:shells-dominated-by-first-moments}
Fix $0\le k\le J-1$ and $u\in I_k$. Set
\[
T:=u-2^k+1
\qquad (1\le T\le 2^k).
\]
Then
\begin{equation}
\label{eq:shells-dominated-by-first-moments}
\abs{Q_{h,k}(u;X)}
\le
\sum_{X\le x\le 2X-T}\abs{B_h(x;T)}+O_h(T^2).
\end{equation}
\end{proposition}

\begin{proof}
Expand the definition of $Q_{h,k}(u;X)$:
\[
Q_{h,k}(u;X)=\sum_{t=2^k}^{u}\sum_{X<n\le 2X-t}a_h(n)a_h(n+t).
\]
Interchanging the order of summation gives
\[
Q_{h,k}(u;X)=\sum_{X<n\le 2X-2^k} a_h(n)\sum_{2^k\le t\le \min(u,2X-n)}a_h(n+t).
\]
Split at $n\le 2X-u$. The main range contributes
\[
\sum_{X<n\le 2X-u} a_h(n)\sum_{t=2^k}^{u}a_h(n+t),
\]
and the inner sum is exactly $B_h(n+2^k-1;T)$. Taking absolute values yields
\[
\abs{\sum_{X<n\le 2X-u} a_h(n)\sum_{t=2^k}^{u}a_h(n+t)}
\le
\sum_{X\le x\le 2X-T}\abs{B_h(x;T)}.
\]
The boundary range $2X-u<n\le 2X-2^k$ has at most $T-1$ values of $n$, and for each such $n$ the inner sum has length at most $T$, so its total contribution is $O(T^2)$. This proves \eqref{eq:shells-dominated-by-first-moments}.
\end{proof}

\begin{theorem}[Short-interval first-moment theorem]
\label{thm:first-moment-criterion}
Assume that
\begin{equation}
\label{eq:first-moment-criterion-hyp}
\sup_{1\le M\le H}
\frac{1}{XM}
\sum_{X\le x\le 2X-M}\abs{B_h(x;M)}\to 0
\qquad (X\to\infty).
\end{equation}
Then \eqref{eq:short-block-criterion-conc} holds.
\end{theorem}

\begin{proof}
Apply \eqref{eq:first-moment-criterion-hyp} at $M=H$. Since $\abs{a_h(n)}\le 1$, every $H$-block satisfies $\abs{B_h(x;H)}\le H$. Hence
\[
\sum_{X\le x\le 2X-H}\abs{B_h(x;H)}^2
\le
H\sum_{X\le x\le 2X-H}\abs{B_h(x;H)}=o(XH^2).
\]
The critical short-block criterion, Theorem~\ref{thm:short-block-criterion}, now gives \eqref{eq:short-block-criterion-conc}.
\end{proof}

\begin{corollary}[First moments also control dyadic local quartic shells]
\label{cor:first-moment-dyadic-shells}
Under the hypothesis \eqref{eq:first-moment-criterion-hyp}, the dyadic local quartic condition \eqref{eq:dyadic-local-quartic-hyp} holds.
\end{corollary}

\begin{proof}
Fix a shell $I_k$ and a point $u\in I_k$, and let
\[
T:=u-2^k+1
\qquad (1\le T\le 2^k\le H).
\]
By Proposition~\ref{prop:shells-dominated-by-first-moments},
\[
\abs{Q_{h,k}(u;X)}
\le
\sum_{X\le x\le 2X-T}\abs{B_h(x;T)}+O_h(T^2).
\]
The hypothesis \eqref{eq:first-moment-criterion-hyp} applies uniformly for every $1\le T\le H$, so
\[
\sum_{X\le x\le 2X-T}\abs{B_h(x;T)}=o(XT).
\]
Since $T\le H\asymp X^{1/2}$, the boundary term also satisfies
\[
T^2\le HT=o(XT).
\]
Hence
\[
\abs{Q_{h,k}(u;X)}=o(XT)=o(X2^k)
\]
uniformly in $k$ and $u$, which is exactly \eqref{eq:dyadic-local-quartic-hyp}.
\end{proof}

\section{A fixed-shift local \texorpdfstring{$U^2$}{U2} reduction}
\label{sec:local-u2}

We next record the local $U^2$ reduction in the normalization used in the formal development. The point which matters for the sequel is that the normalized local quantity below is a fourth-power quantity. Consequently the correct block estimate is a fourth-power estimate, not a linear estimate in the normalized fourth-power mass.

For a finitely supported function $f:\Z\to\C$ and a finite interval $I\subset\Z$, put
\[
C_f(I;t):=\sum_{n\in I}\1_I(n+t)f(n+t)\overline{f(n)}.
\]
Define
\[
U_2^4(f;I):=\sum_{-|I|\le t\le |I|}\abs{C_f(I;t)}^2,
\qquad
\mathcal U_2(f;I):=\frac{U_2^4(f;I)}{\max(|I|,1)^3}.
\]
For $x,M\in\N$, let
\[
I_{x,M}:=\{x+1,x+2,\ldots,x+M\}\subset\Z
\]
and let $\widetilde a:\Z\to\C$ be the zero extension of $a$ from $\N$ to $\Z$, viewed as complex-valued. Thus
\[
B_a(x;M)=\sum_{x<n\le x+M}a(n)=\sum_{n\in I_{x,M}}\widetilde a(n).
\]
When $a=a_h$, we write $B_h(x;M)$ as before.

\begin{proposition}[Correlation expansion]
\label{prop:local-u2-correlation-expansion}
For every $a:\N\to\R$ and every $x,M\in\N$,
\[
\sum_{-M\le t\le M} C_{\widetilde a}(I_{x,M};t)
=
\left(\sum_{n\in I_{x,M}}\widetilde a(n)\right)
\overline{\left(\sum_{n\in I_{x,M}}\widetilde a(n)\right)}.
\]
Equivalently,
\[
\sum_{-M\le t\le M} C_{\widetilde a}(I_{x,M};t)
=B_a(x;M)\overline{B_a(x;M)}.
\]
\end{proposition}

\begin{proof}
Expand the left-hand side and interchange the two finite sums:
\[
\sum_{-M\le t\le M}\sum_{n\in I_{x,M}}
\1_{I_{x,M}}(n+t)\widetilde a(n+t)\overline{\widetilde a(n)}.
\]
For fixed $n\in I_{x,M}$, the change of variables $m=n+t$ bijects the admissible shifts with $m\in I_{x,M}$. Hence the inner sum is
\[
\sum_{m\in I_{x,M}}\widetilde a(m)\overline{\widetilde a(n)}.
\]
Summing in $n$ gives
\[
\sum_{n\in I_{x,M}}\left(\sum_{m\in I_{x,M}}\widetilde a(m)\right)
\overline{\widetilde a(n)}
=
\left(\sum_{m\in I_{x,M}}\widetilde a(m)\right)
\overline{\left(\sum_{n\in I_{x,M}}\widetilde a(n)\right)}.
\]
\end{proof}

\begin{proposition}[Fourth-power local $U^2$ block estimate]
\label{prop:local-u2-fourth-power-block}
For every $a:\N\to\R$, every $x\in\N$, and every $M\ge1$,
\[
\abs{B_a(x;M)}^4
\le
3M^4\,\mathcal U_2(\widetilde a;I_{x,M}).
\]
\end{proposition}

\begin{proof}
By Proposition~\ref{prop:local-u2-correlation-expansion},
\[
\abs{B_a(x;M)}^4
=
\left|\sum_{-M\le t\le M}C_{\widetilde a}(I_{x,M};t)\right|^2.
\]
Cauchy's inequality gives
\[
\left|\sum_{-M\le t\le M}C_{\widetilde a}(I_{x,M};t)\right|^2
\le
(2M+1)\sum_{-M\le t\le M}\abs{C_{\widetilde a}(I_{x,M};t)}^2.
\]
Since $M\ge1$, $2M+1\le3M$. Therefore
\[
\abs{B_a(x;M)}^4
\le
3M\,U_2^4(\widetilde a;I_{x,M}).
\]
Finally $|I_{x,M}|=M$, so
\[
U_2^4(\widetilde a;I_{x,M})
=M^3\mathcal U_2(\widetilde a;I_{x,M}),
\]
which gives the claimed bound.
\end{proof}

\begin{remark}[The missing fourth root]
\label{rem:local-u2-fourth-root}
Because $\mathcal U_2$ is normalized as a fourth-power quantity, Proposition~\ref{prop:local-u2-fourth-power-block} implies only
\[
\abs{B_a(x;M)}
\le
3^{1/4}M\,\mathcal U_2(\widetilde a;I_{x,M})^{1/4}.
\]
A pointwise estimate of the form
\[
\abs{B_a(x;M)}\le C M\,\mathcal U_2(\widetilde a;I_{x,M})
\]
would use the fourth-power normalized quantity linearly. That is not the estimate used below. The averaged first-moment implication is instead obtained from the fourth-power estimate by a scalar Young inequality.
\end{remark}

\begin{lemma}[Young form of the fourth-root step]
\label{lem:young-fourth-scaled}
Let $y,q,M,\eta\ge0$, with $M>0$ and $\eta>0$. If
\[
y^4\le3M^4q,
\]
then
\[
y\le \eta M+\frac{3}{\eta^3}Mq.
\]
\end{lemma}

\begin{proof}
Put $t=y/M$. Then $t\ge0$ and $t^4\le3q$. If $t\le\eta$, the claim follows immediately. Otherwise $\eta\le t$, hence
\[
t\le \frac{t^4}{\eta^3}\le\frac{3q}{\eta^3}.
\]
Multiplying by $M$ and adding the harmless term $\eta M$ gives the displayed estimate in both cases.
\end{proof}

\begin{theorem}[Uniform fourth-power local $U^2$ criterion for fixed $h$]
\label{thm:local-u2-criterion}
Assume that
\begin{equation}
\label{eq:local-u2-criterion-hyp}
\sup_{1\le M\le H}
\frac{1}{X}
\sum_{X\le x\le 2X-M}
\mathcal U_2(\widetilde a_h;I_{x,M})
\to0
\qquad (X\to\infty).
\end{equation}
Then
\[
A_h(N)=o(N)
\qquad (N\to\infty).
\]
\end{theorem}

\begin{proof}
We first derive the first-moment hypothesis. Fix $\varepsilon>0$, set $\eta:=\varepsilon/2$, and put
\[
K:=\frac{3}{\eta^3}.
\]
Apply \eqref{eq:local-u2-criterion-hyp} with the tolerance $\delta:=\varepsilon/(2K)$. For all sufficiently large $X$, uniformly for $1\le M\le H$,
\[
\sum_{X\le x\le2X-M}\mathcal U_2(\widetilde a_h;I_{x,M})
\le \delta X.
\]
For each such $x,M$, Proposition~\ref{prop:local-u2-fourth-power-block} and Lemma~\ref{lem:young-fourth-scaled} give
\[
\abs{B_h(x;M)}
\le
\eta M+K M\mathcal U_2(\widetilde a_h;I_{x,M}).
\]
Summing over $X\le x\le2X-M$, and using that this interval has at most $X$ integer starting points, gives
\[
\sum_{X\le x\le2X-M}\abs{B_h(x;M)}
\le
\eta MX+KM\sum_{X\le x\le2X-M}\mathcal U_2(\widetilde a_h;I_{x,M}).
\]
The last sum is at most $\delta X$, hence
\[
\sum_{X\le x\le2X-M}\abs{B_h(x;M)}
\le
\frac{\varepsilon}{2}MX+KM\frac{\varepsilon}{2K}X
=\varepsilon XM.
\]
This is precisely \eqref{eq:first-moment-criterion-hyp}. The conclusion follows from Theorem~\ref{thm:first-moment-criterion}.
\end{proof}

\begin{remark}[Why the natural $U^3$ upgrade still misses fixed $h$]
A tempting way to force the local input is to rewrite the block sum as a mean of a discrete derivative and then try to control that derivative by a local $U^3$ norm of $\lambda$. On a finite abelian group $G$ one has the exact identity from Gowers-norm formalism \cite{gowers2001szemeredi,taovu2006additive}
\begin{equation}
\label{eq:u3-derivative-identity}
\|f\|_{U^3(G)}^8 = \mathbb E_{r\in G}\,\|\Delta_r f\|_{U^2(G)}^4,
\qquad
\Delta_r f(x):=f(x+r)\overline{f(x)}.
\end{equation}
Thus a local $U^3$ bound for $\lambda$ controls fourth powers of local $U^2$ quantities for the derivative sequences $\lambda(n)\lambda(n+r)$ only \emph{after averaging over the shift parameter $r$}. In schematic form this yields bounds of the shape
\begin{equation}
\label{eq:schematic-averaged-r-u2}
\frac1X\sum_x\frac1M\sum_{|r|\ll M}
\|\lambda(\cdot)\lambda(\cdot+r)\|_{U^2([x+1,x+M])}^4=o(1),
\end{equation}
which are consistent with a single exceptional fixed shift $r=h$. Consequently the local $U^3$ route proves only an averaged-in-$r$ version of the fixed-shift local $U^2$ criterion, and does not remove the need for a genuinely fixed-shift argument.
\end{remark}

\section{Bad blocks, projected packets, and the remaining bottleneck}
\label{sec:projected-packets}

We now turn to the endpoint of the reduction developed above. Before introducing packets, we record the precise reason why the prime-shear symmetry can only hold after projection to the divisible slice.

For a finite interval $J\subset \Z$, define the local Fourier sum
\[
A_J^{(\alpha)}(\xi):=\sum_{n\in J} g_\alpha(n)\,\ee{-\xi n}.
\]
If $p\nmid h$ is prime, decompose by residue classes modulo $p$:
\[
A_J^{(\alpha)}(\xi)=\sum_{r=0}^{p-1}A_{J,r}^{(\alpha)}(\xi),
\qquad
A_{J,r}^{(\alpha)}(\xi):=\sum_{\substack{n\in J\\ n\equiv r\, (p)}} g_\alpha(n)\,\ee{-\xi n}.
\]
The next proposition is the exact correction to the naive full-block affine-shear statement.

\begin{proposition}[No direct prime-shear consistency on the full block]
\label{prop:no-full-block-affine-shear}
Let $p\nmid h$ be prime and let $J\subset\Z$ be a finite interval. Then the $r=0$ slice satisfies the exact identity
\begin{equation}
\label{eq:projected-local-fourier-shear}
A_{J,0}^{(\alpha)}(\xi)
=
-\,A_{J/p}^{(p^2\alpha)}\!\left(p\xi-\frac{\alpha p(p-1)}{2}\right)
+O(1),
\end{equation}
where $J/p:=\{m\in\Z:\ pm\in J\}$ and the $O(1)$ term is a boundary error. By contrast, for each $r\neq 0$ one has only the tautological expression
\[
A_{J,r}^{(\alpha)}(\xi)
=
\sum_{\substack{m\in\Z\\ pm+r\in J}}
\lambda(pm+r)
\,\ee{Q_\alpha(pm+r)-\xi(pm+r)},
\]
for which complete multiplicativity gives no simplification. In particular, there is no full-block affine-shear law of the shape
\[
A_J^{(\alpha)}(\xi)
\approx
c_{p,\alpha,\xi}\,A_{J'}^{(\beta)}(T_p\xi)
\]
up to boundary error alone, because the contribution of the nonzero residue classes can be as large as the whole sum.
\end{proposition}

\begin{proof}
On the divisible slice, write $n=pm$. Since $\lambda(pm)=-\lambda(m)$, we have
\[
A_{J,0}^{(\alpha)}(\xi)
=
-\sum_{m\in J/p}\lambda(m)\,\ee{Q_\alpha(pm)-p\xi m} + O(1),
\]
where the $O(1)$ term accounts for endpoints when $J$ is not exactly a union of residue classes modulo $p$. A direct calculation gives
\[
Q_\alpha(pm)-Q_{p^2\alpha}(m)=\frac{\alpha p(p-1)}{2}\,m,
\]
whence
\[
Q_\alpha(pm)-p\xi m
=
Q_{p^2\alpha}(m)-\left(p\xi-\frac{\alpha p(p-1)}{2}\right)m.
\]
Substituting this identity proves \eqref{eq:projected-local-fourier-shear}.

For $r\neq 0$, writing $n=pm+r$ yields
\[
A_{J,r}^{(\alpha)}(\xi)
=
\sum_{\substack{m\in\Z\\ pm+r\in J}}
\lambda(pm+r)
\,\ee{Q_\alpha(pm+r)-\xi(pm+r)},
\]
and there is no relation between $\lambda(pm+r)$ and $\lambda(m)$ furnished by complete multiplicativity. Therefore the uncontrolled nonzero residue classes cannot be absorbed into a boundary error, and the full-block sum admits no exact analogue of the divisible-slice shear identity.
\end{proof}

\begin{remark}[What the preceding proposition means]
Proposition~\ref{prop:no-full-block-affine-shear} shows that the arithmetic self-similarity lives only on the sparse slice $n\equiv 0 \pmod p$. For the original derivative sequence $a_h(n)=\lambda(n)\lambda(n+h)$ the situation is even less symmetric, since on that slice one has
\[
a_h(pm)=\lambda(pm)\lambda(pm+h),
\]
which does not simplify multiplicatively when $p\nmid h$. Thus the correct next target is not a full-block prime-shear law, but a theorem that manufactures a projector or weight selecting enough of the divisible slice while preserving the mass coming from a bad fixed-shift block.
\end{remark}

Let $I\subset\Z$ be a finite interval of length $L$, let $\alpha\in\R$, and define
\[
S_I(\alpha):=\sum_{n\in I}a_h(n)\,\ee{\alpha n}.
\]
To linearize the shift $h$, introduce the quadratic phase
\[
Q_\alpha(n):=\frac{\alpha}{2h}n(n-h)
\qquad (n\in\Z),
\]
so that
\[
Q_\alpha(n+h)-Q_\alpha(n)=\alpha n.
\]
Define
\[
g_\alpha(n):=\lambda(n)\,\ee{Q_\alpha(n)}.
\]
Then
\[
g_\alpha(n+h)\overline{g_\alpha(n)}=a_h(n)\,\ee{\alpha n},
\]
and therefore
\[
S_I(\alpha)=\sum_{n\in I}g_\alpha(n+h)\overline{g_\alpha(n)}.
\]

\begin{proposition}[A bad block forces an unprojected packet]
\label{prop:bad-block-packet}
Define
\[
\mathcal D_{I,\alpha}
:=
\sum_{|r|<L}(L-|r|)\,\ee{\alpha r}
\sum_{n\in\Z}a_h(n+r)a_h(n)\,\mathbf 1_I(n+r)\mathbf 1_I(n).
\]
Then
\[
\mathcal D_{I,\alpha}=|S_I(\alpha)|^2.
\]
In particular, if $|S_I(\alpha)|\ge \eta L$, then $\mathcal D_{I,\alpha}\ge \eta^2L^2$.
\end{proposition}

\begin{proof}
Expanding the square gives
\[
|S_I(\alpha)|^2
=
\sum_{u,v\in I}a_h(u)a_h(v)\,\ee{\alpha(u-v)}.
\]
Write $r=u-v$. For fixed $r$ with $|r|<L$, the number of pairs $(u,v)\in I^2$ with $u-v=r$ is $L-|r|$, and there are no such pairs when $|r|\ge L$. Hence
\[
|S_I(\alpha)|^2
=
\sum_{|r|<L}(L-|r|)\,\ee{\alpha r}
\sum_{n\in\Z}a_h(n+r)a_h(n)\,\mathbf 1_I(n+r)\mathbf 1_I(n),
\]
which is exactly $\mathcal D_{I,\alpha}$.
\end{proof}

The packet $\mathcal D_{I,\alpha}$ is the exact unprojected delta/off-diagonal object forced by a bad block.
The next theorem shows that one can already manufacture a genuine projected \emph{delta} packet on a small prime-divisible slice.

\begin{theorem}[Small-prime selector for a projected delta packet]
\label{thm:small-prime-projected-delta}
Fix $0<\eta\le 1$. There exists a function $P_\eta(L)\to\infty$ with $P_\eta(L)=L^{o(1)}$ such that the following holds for all sufficiently large $L$. Let $I\subset\Z$ be an interval of length $L$, let $\alpha\in\R$, and assume that
\[
|S_I(\alpha)|\ge \eta L.
\]
Then there exists a prime $p\le P_\eta(L)$ with $p\nmid h$ such that
\begin{equation}
\label{eq:small-prime-slice-bias}
\left|\sum_{\substack{n\in I\\ p\mid n}} a_h(n)\,\ee{\alpha n}\right|
\gg_\eta \frac{L}{p}.
\end{equation}
Consequently the projected delta packet
\begin{equation}
\label{eq:projected-delta-packet}
\mathcal D^{(p)}_{I,\alpha}
:=
\sum_{r\in\Z}\sum_{n\in\Z}
 a_h(n+r)a_h(n)\,\ee{\alpha r}
\mathbf 1_I(n+r)\mathbf 1_I(n)
\mathbf 1_{p\mid n}\mathbf 1_{p\mid n+r}
\end{equation}
has size
\begin{equation}
\label{eq:projected-delta-packet-lb}
\mathcal D^{(p)}_{I,\alpha}
\gg_\eta
\frac{L^2}{p^2}.
\end{equation}
Since $p\mid n$ and $p\mid n+r$ force $r=p\ell$, one may also rewrite \eqref{eq:projected-delta-packet} as
\begin{equation}
\label{eq:projected-delta-packet-step-p}
\mathcal D^{(p)}_{I,\alpha}
=
\sum_{\ell\in\Z}\ee{\alpha p\ell}
\sum_{n\in\Z} a_h(n+p\ell)a_h(n)
\mathbf 1_I(n+p\ell)\mathbf 1_I(n)\mathbf 1_{p\mid n}.
\end{equation}
\end{theorem}

\begin{proof}
Choose any function $P_\eta(L)\to\infty$ satisfying
\[
\sum_{\substack{p\le P_\eta(L)\\ p\nmid h}}\frac1p\to\infty
\qquad\text{and}\qquad
\frac{\pi(P_\eta(L))^2}{L}\to 0
\qquad (L\to\infty),
\]
for instance $P_\eta(L):=\exp(\sqrt{\log L})$. Let
\[
\mathcal P:=\{p\le P_\eta(L):\ p\nmid h\},
\qquad
w(n):=\sum_{p\in\mathcal P}\mathbf 1_{p\mid n}.
\]
Write
\[
S:=\sum_{n\in I} a_h(n)\,\ee{\alpha n}=S_I(\alpha),
\qquad
S_p:=\sum_{\substack{n\in I\\ p\mid n}} a_h(n)\,\ee{\alpha n},
\qquad
\mu:=\frac1L\sum_{n\in I}w(n).
\]
Then
\[
\sum_{p\in\mathcal P}S_p=\sum_{n\in I} w(n)a_h(n)\,\ee{\alpha n}.
\]
Centering $w$ at its mean gives
\[
\sum_{p\in\mathcal P}S_p
=
\mu S+
\sum_{n\in I}(w(n)-\mu)a_h(n)\,\ee{\alpha n}.
\]
Hence, by Cauchy--Schwarz,
\[
\left|\sum_{p\in\mathcal P}S_p\right|
\ge
\mu|S|-
\left(\sum_{n\in I}(w(n)-\mu)^2\right)^{1/2}L^{1/2}
=
\mu|S|-\sigma L,
\]
where
\[
\sigma^2:=\frac1L\sum_{n\in I}(w(n)-\mu)^2.
\]
The standard first and second moment estimates for the divisor-counting weight $w$ give
\[
\mu=
\sum_{p\in\mathcal P}\frac1p+O\!\left(\frac{\pi(P_\eta(L))}{L}\right)
\]
and
\[
\sigma^2\ll
\sum_{p\in\mathcal P}\frac1p
+
\frac{\pi(P_\eta(L))^2}{L}.
\]
Because $\sum_{p\in\mathcal P}1/p\to\infty$ and $\pi(P_\eta(L))^2/L\to 0$, we have $\sigma/\mu\to 0$. Therefore, for all sufficiently large $L$,
\[
\left|\sum_{p\in\mathcal P}S_p\right|\ge \frac{\eta}{2}\mu L.
\]
If every $p\in\mathcal P$ satisfied $|S_p|<(\eta/2)N_p(I)$, where
\[
N_p(I):=\#\{n\in I:\ p\mid n\},
\]
then
\[
\sum_{p\in\mathcal P}|S_p|<\frac{\eta}{2}\sum_{p\in\mathcal P}N_p(I)=\frac{\eta}{2}\mu L,
\]
contradicting the previous lower bound and the triangle inequality. Hence there exists $p\in\mathcal P$ such that
\[
|S_p|\ge \frac{\eta}{2}N_p(I)\gg_\eta \frac{L}{p},
\]
which proves \eqref{eq:small-prime-slice-bias}.

Finally,
\[
|S_p|^2
=
\sum_{u,v\in\Z}
a_h(u)a_h(v)\,\ee{\alpha(u-v)}
\mathbf 1_I(u)\mathbf 1_I(v)
\mathbf 1_{p\mid u}\mathbf 1_{p\mid v}.
\]
Writing $u=n+r$ and $v=n$ yields exactly \eqref{eq:projected-delta-packet}. Therefore
\[
\mathcal D^{(p)}_{I,\alpha}=|S_p|^2\gg_\eta \frac{L^2}{p^2},
\]
which is \eqref{eq:projected-delta-packet-lb}. The representation \eqref{eq:projected-delta-packet-step-p} follows because $p\mid n$ and $p\mid n+r$ imply $p\mid r$.
\end{proof}

\begin{remark}[What the small-prime selector does and does not give]
Theorem~\ref{thm:small-prime-projected-delta} produces a genuine projected delta/off-diagonal object from a bad block. However, its kernel is the bare correlation
\[
a_h(n+p\ell)a_h(n)\,\ee{\alpha p\ell},
\]
whereas the exact affine-shear mechanism acts on projected packets built from
\[
g_\alpha(n+p\ell h)\overline{g_\alpha(n)}.
\]
Thus the theorem does not yet produce the projected Fej\'er/shear packet required for the quadratic-obstruction argument below. It closes the ``manufacture a projector'' step only at the delta-packet level, not at the Fej\'er/shear level.
\end{remark}

To bring in the exact arithmetic self-similarity that survives after projection, we now work with the projected Fej\'er/shear packet on the divisible slice modulo a prime $p\nmid h$.

\begin{proposition}[Exact affine shear on the divisible slice]
\label{prop:projected-affine-shear}
Let $p\nmid h$ be prime, let $R\ge 1$ be an integer, and let $z\in\C$ with $|z|=1$. Define
\[
\mathcal F^{(p)}_{I,\alpha}(R;z)
:=
\sum_{|m|<pR}\left(1-\frac{|m|}{pR}\right)z^m
\sum_{n\in\Z}
\mathbf 1_I(n)\mathbf 1_I(n+mh)
\mathbf 1_{p\mid n}\mathbf 1_{p\mid n+mh}
 g_\alpha(n+mh)\overline{g_\alpha(n)}.
\]
Also define
\[
\beta_p:=\frac{\alpha p(p-1)}{2},
\qquad
I/p:=\{n'\in\Z:\ pn'\in I\}.
\]
Then
\[
\mathcal F^{(p)}_{I,\alpha}(R;z)
=
\sum_{|\ell|<R}\left(1-\frac{|\ell|}{R}\right)
\bigl(z^p\ee{\beta_ph}\bigr)^\ell
\sum_{n'\in\Z}
\mathbf 1_{I/p}(n')\mathbf 1_{I/p}(n'+\ell h)
 g_{p^2\alpha}(n'+\ell h)\overline{g_{p^2\alpha}(n')}.
\]
\end{proposition}

\begin{proof}
Because $p\nmid h$, the conditions $p\mid n$ and $p\mid n+mh$ imply $p\mid m$. Write $m=p\ell$ and $n=pn'$. By complete multiplicativity of $\lambda$,
\[
g_\alpha(pm)
=
-\lambda(m)\,\ee{Q_\alpha(pm)}.
\]
Also
\[
Q_\alpha(pm)=Q_{p^2\alpha}(m)+\beta_p m,
\]
so
\[
g_\alpha(pm)=-\ee{\beta_p m}\,g_{p^2\alpha}(m).
\]
Substituting $m=p\ell$ and $n=pn'$ into the definition of $\mathcal F^{(p)}_{I,\alpha}(R;z)$, the two minus signs cancel and we obtain
\[
g_\alpha(p(n'+\ell h))\overline{g_\alpha(pn')}
=
\ee{\beta_p\ell h}\,g_{p^2\alpha}(n'+\ell h)\overline{g_{p^2\alpha}(n')}.
\]
The stated identity follows.
\end{proof}


\subsection{Projected transport defect and rigidity}

Fix a prime $p\nmid h$, a dyadic interval $J_0\subset\Z$ of length $|J_0|=2^L$, and a phase $\alpha\in\T$. Put
\[
\beta_p:=\frac{\alpha p(p-1)}{2},
\qquad
b(n):=\ee{-\beta_ph}\,g_\alpha(p(n+h))\overline{g_\alpha(pn)}
= g_{p^2\alpha}(n+h)\overline{g_{p^2\alpha}(n)}.
\]
Then $|b(n)|=1$ for every $n$. For each dyadic subinterval $J\subseteq J_0$, define
\[
m_J:=\frac1{|J|}\sum_{n\in J} b(n),
\qquad
c_J:=\frac{m_J}{|m_J|}\in\T
\quad (m_J\neq 0).
\]
If $J$ has dyadic children $J_-,J_+$, define the transport defect
\[
\delta(J):=\frac{|m_{J_-}|+|m_{J_+}|}{2}-|m_J|\ge 0.
\]
Write $\mathcal D^\circ(J_0)$ for the set of proper dyadic subintervals of $J_0$, and define the total transport defect
\[
\mathsf T(J_0):=\frac1{|J_0|}\sum_{J\in\mathcal D^\circ(J_0)} |J|\,\delta(J).
\]
Finally, for $m\ge 0$ let
\[
J_0(m):=\{n\in J_0:\ n,n+h,\dots,n+mh\in J_0\}.
\]

\begin{proposition}[Small projected transport defect implies projected rigidity]
\label{prop:projected-rigidity}
Assume that $|m_J|\ge \eta$ for every dyadic subinterval $J\subseteq J_0$, where $0<\eta\le 1$. Then with $c:=c_{J_0}$ one has
\[
\frac1{|J_0|}\sum_{n\in J_0}|b(n)-c|
\le
\varepsilon,
\qquad
\varepsilon:=\frac{\sqrt{8L\,\mathsf T(J_0)}}{\eta}.
\]
\end{proposition}

\begin{proof}
Let $J\subseteq J_0$ be dyadic with children $J_-,J_+$. Write
\[
m_{J_\pm}=a_\pm u_\pm,
\qquad
a_\pm:=|m_{J_\pm}|\in [\eta,1],
\qquad
u_\pm\in\T,
\]
and write $m_J=|m_J|c_J$ with $c_J\in\T$. Since
\[
m_J=\frac{m_{J_-}+m_{J_+}}{2},
\]
we have
\[
4|m_J|^2=(a_-+a_+)^2-a_-a_+|u_--u_+|^2.
\]
Also,
\[
2\delta(J)=a_-+a_+-2|m_J|.
\]
Hence
\[
2\delta(J)=\frac{a_-a_+|u_--u_+|^2}{a_-+a_++2|m_J|}
\ge \frac{\eta^2}{4}|u_--u_+|^2,
\]
so
\[
|u_--u_+|^2\le \frac{8}{\eta^2}\,\delta(J).
\]
Because $c_J$ is the phase of the convex combination $a_-u_-+a_+u_+$, it lies on the shorter arc joining $u_-$ and $u_+$. Therefore
\[
|c_{J_\pm}-c_J|\le |u_--u_+|,
\qquad
|c_{J_\pm}-c_J|^2\le \frac{8}{\eta^2}\,\delta(J).
\]
Now fix $n\in J_0$, and let
\[
J_0\supset J_1(n)\supset \cdots \supset J_L(n)=\{n\}
\]
be the dyadic ancestor chain of $n$. Since $c_{\{n\}}=b(n)$, telescoping gives
\[
|b(n)-c_{J_0}|
\le \sum_{j=0}^{L-1}|c_{J_{j+1}(n)}-c_{J_j(n)}|.
\]
By Cauchy--Schwarz,
\[
|b(n)-c_{J_0}|^2
\le L\sum_{j=0}^{L-1}|c_{J_{j+1}(n)}-c_{J_j(n)}|^2.
\]
Average over $n\in J_0$. Each internal dyadic interval $J\subsetneq J_0$ contributes with weight $|J|/|J_0|$, so
\[
\frac1{|J_0|}\sum_{n\in J_0}|b(n)-c_{J_0}|^2
\le \frac{8L}{\eta^2|J_0|}
\sum_{J\in\mathcal D^\circ(J_0)} |J|\,\delta(J)
=
\frac{8L}{\eta^2}\,\mathsf T(J_0).
\]
Taking square roots and using Cauchy--Schwarz once more gives the claim.
\end{proof}

\begin{proposition}[Projected rigidity yields a long projected Fej\'er packet]
\label{prop:projected-rigidity-to-packet}
Assume the hypotheses of Proposition~\ref{prop:projected-rigidity}. Let $R\ge 1$ satisfy
\[
R\le \frac{|J_0|}{8|h|}
\qquad\text{and}\qquad
\mathsf T(J_0)\le \frac{\eta^2}{512LR^2}.
\]
Then, with $c:=c_{J_0}$,
\[
\Re\sum_{m=0}^{R-1}\left(1-\frac{m}{R}\right)
\overline c^{\,m}
\frac1{|J_0(m)|}
\sum_{n\in J_0(m)} g_{p^2\alpha}(n+mh)\overline{g_{p^2\alpha}(n)}
\ge \frac{R}{4}.
\]
Consequently, if $I:=pJ_0$ and $z\in\T$ is chosen so that
\[
z^p\ee{\beta_ph}=\overline c,
\]
then
\[
\Re\,\mathcal F^{(p)}_{I,\alpha}(R;z)
\ge \frac{R|J_0|}{8}.
\]
\end{proposition}

\begin{proof}
By Proposition~\ref{prop:projected-rigidity},
\[
\frac1{|J_0|}\sum_{n\in J_0}|b(n)-c|\le \varepsilon
:=\frac{\sqrt{8L\,\mathsf T(J_0)}}{\eta}
\le \frac{1}{8R}.
\]
For $0\le m<R$ and $n\in J_0(m)$ one has the telescoping identity
\[
g_{p^2\alpha}(n+mh)\overline{g_{p^2\alpha}(n)}\,\overline c^{\,m}
=
\prod_{j=0}^{m-1}\bigl(b(n+jh)\overline c\bigr).
\]
Since each factor has modulus at most $1$,
\[
\left|\prod_{j=0}^{m-1} z_j-1\right|\le \sum_{j=0}^{m-1}|z_j-1|
\qquad (|z_j|\le 1).
\]
Therefore
\[
\left|
\frac1{|J_0(m)|}\sum_{n\in J_0(m)}g_{p^2\alpha}(n+mh)\overline{g_{p^2\alpha}(n)}\,\overline c^{\,m}-1
\right|
\le
\sum_{j=0}^{m-1}
\frac1{|J_0(m)|}\sum_{n\in J_0(m)}|b(n+jh)-c|.
\]
Because $m<R\le |J_0|/(8|h|)$, we have $|J_0(m)|\ge |J_0|/2$. Hence each inner average is at most $2\varepsilon$, and so
\[
\left|
\frac1{|J_0(m)|}\sum_{n\in J_0(m)}g_{p^2\alpha}(n+mh)\overline{g_{p^2\alpha}(n)}\,\overline c^{\,m}-1
\right|
\le 2m\varepsilon.
\]
Taking real parts and summing against Fej\'er weights yields
\[
\Re\sum_{m=0}^{R-1}\left(1-\frac{m}{R}\right)
\overline c^{\,m}
\frac1{|J_0(m)|}
\sum_{n\in J_0(m)} g_{p^2\alpha}(n+mh)\overline{g_{p^2\alpha}(n)}
\ge
\sum_{m=0}^{R-1}\left(1-\frac{m}{R}\right)(1-2m\varepsilon).
\]
Since $R\varepsilon\le 1/8$,
\[
\sum_{m=0}^{R-1}\left(1-\frac{m}{R}\right)(1-2m\varepsilon)
\ge \frac{R}{2}-2\varepsilon\sum_{m=0}^{R-1}m
\ge \frac{R}{2}-\varepsilon R^2
\ge \frac{R}{4}.
\]
This proves the first claim.

For the second, note that $I=pJ_0$ and Proposition~\ref{prop:projected-affine-shear} give
\[
\mathcal F^{(p)}_{I,\alpha}(R;z)
=
\sum_{|m|<R}\left(1-\frac{|m|}{R}\right)
\bigl(z^p\ee{\beta_ph}\bigr)^m
\sum_{n\in\Z}\1_{J_0}(n)\1_{J_0}(n+mh)
 g_{p^2\alpha}(n+mh)\overline{g_{p^2\alpha}(n)}.
\]
By the choice of $z$, the phase factor equals $\overline c^{\,m}$. Discarding the negative values of $m$ and using $|J_0(m)|\ge |J_0|/2$ for $0\le m<R$, we obtain
\[
\Re\,\mathcal F^{(p)}_{I,\alpha}(R;z)
\ge \frac{|J_0|}{2}
\Re\sum_{m=0}^{R-1}\left(1-\frac{m}{R}\right)
\overline c^{\,m}
\frac1{|J_0(m)|}
\sum_{n\in J_0(m)} g_{p^2\alpha}(n+mh)\overline{g_{p^2\alpha}(n)}
\ge \frac{R|J_0|}{8}.
\]
This proves the proposition.
\end{proof}

\begin{proposition}[Exact transport-defect identity on a projected slice]
\label{prop:projected-transport-identity}
With the notation above,
\[
\sum_{J\in\mathcal D^\circ(J_0)} |J|\,\delta(J)
=
|J_0|-|J_0|\,|m_{J_0}|.
\]
Equivalently,
\[
\mathsf T(J_0)=1-|m_{J_0}|.
\]
\end{proposition}

\begin{proof}
By definition,
\[
|J|\,\delta(J)
=
\frac{|J|}{2}\bigl(|m_{J_-}|+|m_{J_+}|\bigr)-|J|\,|m_J|
=
|J_-|\,|m_{J_-}|+|J_+|\,|m_{J_+}|-|J|\,|m_J|.
\]
Summing over all proper dyadic subintervals $J\subsetneq J_0$ telescopes to
\[
\sum_{J\in\mathcal D^\circ(J_0)} |J|\,\delta(J)
=
\sum_{\text{dyadic leaves }\{n\}\subset J_0} |\{n\}|\,|m_{\{n\}}|-|J_0|\,|m_{J_0}|.
\]
But $m_{\{n\}}=b(n)$ and $|b(n)|=1$, so the leaf sum equals $|J_0|$. This gives the identity.
\end{proof}

\begin{corollary}[Greedy chain selection from moderate projected bias]
\label{cor:projected-greedy-chain}
Let $\rho:=|m_{J_0}|$, and build a greedy chain
\[
J_0\supset J_1\supset J_2\supset \cdots
\]
by always choosing $J_{k+1}$ to be the child of $J_k$ with larger $|m|$. Then for every $r\ge 1$ one has
\[
\sum_{k=0}^{r-1}\delta(J_k)\le 1-\rho.
\]
Consequently, for any threshold $\lambda>0$,
\[
\#\{0\le k<r:\ \delta(J_k)>\lambda\}
\le \frac{1-\rho}{\lambda}.
\]
In particular, among the first $r$ levels there is a consecutive run of length at least
\[
\frac{r}{\lfloor (1-\rho)/\lambda\rfloor+1}-1
\]
on which every defect is at most $\lambda$.
\end{corollary}

\begin{proof}
For each $k$,
\[
|m_{J_{k+1}}|
\ge \frac{|m_{(J_k)_-}|+|m_{(J_k)_+}|}{2}
= |m_{J_k}|+\delta(J_k).
\]
Summing this from $k=0$ to $r-1$ and using $|m_{J_r}|\le 1$ gives
\[
\sum_{k=0}^{r-1}\delta(J_k)
\le 1-|m_{J_0}|=1-\rho.
\]
The bound on the number of indices with $\delta(J_k)>\lambda$ is immediate. If there are at most $q:=\lfloor (1-\rho)/\lambda\rfloor$ bad indices among the first $r$ levels, then these indices split the range $\{0,1,\dots,r-1\}$ into at most $q+1$ consecutive blocks, so one block has length at least $r/(q+1)-1$.
\end{proof}

\begin{proposition}[A large projected packet forces a quadratic obstruction]
\label{prop:projected-packet-quadratic}
Let $J:=I/p$. Assume that $R\asymp |J|/|h|$ and that for some $\kappa>0$ and some unit complex number $z$ one has
\[
\Re\,\mathcal F^{(p)}_{I,\alpha}(R;z)\ge \kappa R|J|.
\]
Then
\[
\sup_{\deg P\le 2}
\left|\frac{1}{|J|}\sum_{n\in J}\lambda(n)\,\ee{P(n)}\right|
\gg_{\kappa,h}1.
\]
\end{proposition}

\begin{proof}
Choose $\gamma\in\R$ such that
\[
\ee{\gamma h}=z^p\ee{\beta_p h}.
\]
Define
\[
f(n):=g_{p^2\alpha}(n)\mathbf 1_J(n),
\qquad
\widetilde f(n):=f(n)\ee{\gamma n}.
\]
By Proposition~\ref{prop:projected-affine-shear},
\[
\mathcal F^{(p)}_{I,\alpha}(R;z)
=
\sum_{|\ell|<R}\left(1-\frac{|\ell|}{R}\right)
\sum_{n\in\Z}\widetilde f(n+\ell h)\overline{\widetilde f(n)}.
\]
Embed $J$ into a cyclic group $G=\Z/N\Z$ with $N\asymp |J|$ and no wraparound. Let
\[
\widehat{\widetilde f}(\xi):=\frac1N\sum_{n\in G}\widetilde f(n)\,\ee{-\xi n/N}
\qquad (\xi\in G).
\]
Fourier inversion gives
\[
\sum_{|\ell|<R}\left(1-\frac{|\ell|}{R}\right)
\sum_{n\in G}\widetilde f(n+\ell h)\overline{\widetilde f(n)}
=
N\sum_{\xi\in G}|\widehat{\widetilde f}(\xi)|^2F_R\!\left(\frac{h\xi}{N}\right),
\]
where
\[
F_R(t):=\sum_{|\ell|<R}\left(1-\frac{|\ell|}{R}\right)\ee{\ell t}
\]
is the Fej\'er kernel. Hence
\[
N\sum_{\xi\in G}|\widehat{\widetilde f}(\xi)|^2F_R\!\left(\frac{h\xi}{N}\right)
\ge \kappa R|J|.
\]
Since $\sum_{\xi\in G}|\widehat{\widetilde f}(\xi)|^2=|J|/N$ by Parseval and $F_R\ge 0$, there is a constant $c_h>0$ such that the set
\[
\Omega_h:=\left\{\xi\in G: F_R\!\left(\frac{h\xi}{N}\right)\ge c_hR\right\}
\]
has cardinality $O_h(1)$ when $R\asymp N/|h|$. Therefore some $\xi_0\in\Omega_h$ satisfies
\[
|\widehat{\widetilde f}(\xi_0)|^2\gg_{\kappa,h}1.
\]
Unwinding the definitions,
\[
\widehat{\widetilde f}(\xi_0)
=
\frac1N\sum_{n\in J}\lambda(n)
\,\ee{Q_{p^2\alpha}(n)+\gamma n-\xi_0n/N}.
\]
The phase
\[
P(n):=Q_{p^2\alpha}(n)+\gamma n-\xi_0n/N
\]
has degree at most $2$, and $N\asymp |J|$, so
\[
\sup_{\deg P\le 2}
\left|\frac{1}{|J|}\sum_{n\in J}\lambda(n)\,\ee{P(n)}\right|
\gg_{\kappa,h}1.
\]
This proves the claim.
\end{proof}

This conclusion should be viewed as a local quadratic-phase obstruction in the spirit of the polynomial-phase and nilsequence literature \cite{greentao2012mobiusnil,matomaki2023higher,matomakiShao2021poly}.

\begin{remark}[Remaining bottleneck]
\label{rem:current-bottleneck}
The rigorous chain established above isolates the remaining obstruction. Indeed,
\[
\text{bad block}
\Longrightarrow
\text{large unprojected packet}
\]
by Proposition~\ref{prop:bad-block-packet}, and also
\[
\text{bad block}
\Longrightarrow
\text{large projected delta packet on }0\!\!\pmod p
\]
for some small prime $p\nmid h$, by Theorem~\ref{thm:small-prime-projected-delta}. On the projected side we now also have the deterministic implication
\begin{align*}
&\text{tiny projected transport defect}
\Longrightarrow
\text{projected rigidity}\\
&\Longrightarrow
\text{large projected Fej\'er/shear packet}
\Longrightarrow
\text{local quadratic obstruction},
\end{align*}
by Propositions~\ref{prop:projected-rigidity}, \ref{prop:projected-rigidity-to-packet}, and \ref{prop:projected-packet-quadratic}.

However, Proposition~\ref{prop:projected-transport-identity} shows that tiny total projected transport defect is equivalent to near-total projected top mean:
\[
\mathsf T(J_0)=1-|m_{J_0}|.
\]
Thus a selector that only produces a projected slice with $|m_{J_0}|\gg_\eta 1$ cannot by itself yield the transport bound needed for Proposition~\ref{prop:projected-rigidity-to-packet}. The strongest unconditional replacement currently available is Corollary~\ref{cor:projected-greedy-chain}: from moderate projected bias one can still extract a long run of small \emph{one-step} defects along a greedy chain.

The remaining gap is therefore sharper than the earlier ``delta packet to Fej\'er packet'' bridge. One must either:
\begin{enumerate}[label=(\roman*)]
\item manufacture a projector or weight that boosts the projected top mean from $\gg_\eta 1$ to $1-o(1)$, so that Proposition~\ref{prop:projected-rigidity-to-packet} applies at the full length $R\asymp |J_0|/|h|$; or
\item prove a stronger contagion theorem that upgrades a long run of small one-step defects, as in Corollary~\ref{cor:projected-greedy-chain}, directly into a long projected Fej\'er packet.
\end{enumerate}
This is the remaining bottleneck in the present argument.
\end{remark}

\section{Proof of the conditional fixed-shift theorem}

\begin{proof}[Proof of Theorem~\ref{thm:main}]
The hypothesis of Theorem~\ref{thm:main} is precisely the short-interval first-moment assumption \eqref{eq:first-moment-criterion-hyp}. By Theorem~\ref{thm:first-moment-criterion}, this yields \eqref{eq:short-block-criterion-conc}. Theorem~\ref{thm:short-block-criterion} then applies directly and gives
\[
A_h(N)=o(N)
\qquad (N\to\infty).
\]
\end{proof}

\begin{remark}
There are two routes through the paper. The direct route proving Theorem~\ref{thm:main} is
\begin{center}
absolute first moment
$\Longrightarrow$ critical short-block energy
$\Longrightarrow$ prefix cancellation.
\end{center}
The structural route is
\begin{center}
short-block criterion
$\Longrightarrow$ energy profile and discrete derivative
$\Longrightarrow$ Fourier--Dirichlet reformulation
$\Longrightarrow$ positive multiscale estimate
$\Longrightarrow$ dyadic local quartic shells
$\Longrightarrow$ fourth-power local $U^2$ criterion
$\Longrightarrow$ packet reformulations and projected bottleneck.
\end{center}
Theorem~\ref{thm:main} is conditional.\par
The paper ends with the projected-packet bottleneck rather than with an unconditional proof of fixed-shift pair Chowla.
\end{remark}

\section*{Funding}
This research received no external funding.

\section{Acknowledgments}

I would like to thank AGIR Labs for their help and support in my writing this work. Jennifer Dodgson, CEO and Chief AI Engineer, graciously enabled my collaboration with their group. I would also like to thank the DAObi cryptocurrency and Chinese classics project for their help and aid. 

I would like to thank Professor Yang-Hui He for his inspiring work at the intersection of mathematics and artificial intelligence, which influenced aspects of my approach. I also thank Ernest Petherbridge, whose suggestions for literature provided invaluable foundational context for this project.

\begin{thebibliography}{99}

\bibitem[Tao(2016)]{tao2016logarithmically}
T. Tao,
\newblock The logarithmically averaged Chowla and Elliott conjectures for two-point correlations,
\newblock {\em Forum of Mathematics, Pi} \textbf{4} (2016), e8.

\bibitem[Tao and Ter\"av\"ainen(2019)]{taoteravainen2019almostall}
T. Tao and J. Ter\"av\"ainen,
\newblock The structure of correlations of multiplicative functions at almost all scales, with applications to the Chowla and Elliott conjectures,
\newblock {\em Algebra \& Number Theory} \textbf{13} (2019), no.~9, 2103--2150.

\bibitem[Matom\"aki and Radziwi\l\l(2016)]{matomaki2016multiplicative}
K. Matom\"aki and M. Radziwi\l\l,
\newblock Multiplicative functions in short intervals,
\newblock {\em Annals of Mathematics} \textbf{183} (2016), no.~3, 1015--1056.

\bibitem[Matom\"aki, Radziwi\l\l, and Tao(2020)]{matomaki2020fourier}
K. Matom\"aki, M. Radziwi\l\l, and T. Tao,
\newblock Fourier uniformity of bounded multiplicative functions in short intervals on average,
\newblock {\em Inventiones Mathematicae} \textbf{220} (2020), no.~1, 1--58.

\bibitem[Matom\"aki et al.(2023)]{matomaki2023higher}
K. Matom\"aki, M. Radziwi\l\l, T. Tao, J. Ter\"av\"ainen, and T. Ziegler,
\newblock Higher uniformity of bounded multiplicative functions in short intervals on average,
\newblock {\em Annals of Mathematics} \textbf{197} (2023), no.~2, 739--857.

\bibitem[Chowla(1965)]{chowla1965riemann}
S. Chowla,
\newblock {\em The Riemann Hypothesis and Hilbert's Tenth Problem},
\newblock Gordon and Breach, New York, 1965.

\bibitem[Zygmund(1959)]{zygmund1959trigonometric}
A. Zygmund,
\newblock {\em Trigonometric Series},
\newblock 2nd ed., Vols.~I--II, Cambridge University Press, Cambridge, 1959.

\bibitem[Gowers(2001)]{gowers2001szemeredi}
W. T. Gowers,
\newblock A new proof of Szemer\'edi's theorem,
\newblock {\em Geom. Funct. Anal.} \textbf{11} (2001), no.~3, 465--588.

\bibitem[Tao and Vu(2006)]{taovu2006additive}
T. Tao and V. Vu,
\newblock {\em Additive Combinatorics},
\newblock Cambridge Studies in Advanced Mathematics, vol.~105,
\newblock Cambridge University Press, Cambridge, 2006.

\bibitem[Grafakos(2008)]{grafakos2008classical}
L. Grafakos,
\newblock {\em Classical Fourier Analysis},
\newblock 2nd ed., Graduate Texts in Mathematics, vol.~249,
\newblock Springer, New York, 2008.

\bibitem[Green and Tao(2012)]{greentao2012mobiusnil}
B. Green and T. Tao,
\newblock The M\"obius function is strongly orthogonal to nilsequences,
\newblock {\em Ann. of Math. (2)} \textbf{175} (2012), no.~2, 541--566.

\bibitem[Matom\"aki and Shao(2021)]{matomakiShao2021poly}
K. Matom\"aki and X. Shao,
\newblock Discorrelation between primes in short intervals and polynomial phases,
\newblock {\em Int. Math. Res. Not. IMRN} 2021, no.~16, 12330--12355.

\bibitem[Talagrand(2005)]{talagrand2005generic}
M. Talagrand,
\newblock {\em The Generic Chaining: Upper and Lower Bounds of Stochastic Processes},
\newblock Springer Monographs in Mathematics,
\newblock Springer, Berlin, 2005.
\end{thebibliography}

\end{document}